\chapter{Conclusion}
\label{c:conclusion}

\noindent
This thesis explored whether perceived author identity measurably affects the interpretation of
emotion in ambiguous text, for both humans and large language models, and whether the two respond
in comparable ways. Author attribution was manipulated experimentally while the textual content
was held constant. In a first pass, one emotion was selected per condition and the hypotheses were
tested using that label. A second pass kept all eleven emotion proportions, retested the
hypotheses, and identified changes that the first pass had missed.

\section{Summary of Findings}
\label{conc:summary}

\noindent
Author framing shifts human emotion judgments for author-relevant ambiguous texts, on both persona
sides, and has no statistically detectable effect on unambiguous control texts. The second result
fits H1, but it cannot confirm a zero effect or show that the manipulation acted only on author
information. In the author-independent category one persona contrast is significant and the other
is not; a direct comparison does not establish different effect sizes.

The model ensemble is sensitive to author framing on both persona sides. All seven models have
nonzero observed label-change rates, although separate per-model significance tests were not run. No
prompt-configuration effect was detected, so equivalent effects across formats cannot be claimed.
A previously unanalysed generic-human control separates this into two steps under matched models:
adding an unspecified author line alone moves classifications by about 8\% on both sides, and the
step from there to a specific persona moves them by 16--19\%. These are two separately
measured distances, not additive components of one total effect.

The predicted ordering across ambiguity categories is present in raw shift magnitude, and all
three pairwise contrasts survive correction. Referencing each text against its own conditional
sampling null, however, shows that the shift expected under no effect is itself strongly ordered
by category. The excess above that conditional-null expectation is small enough that none of the
contrasts remains significant. The raw ordering agrees with H3's prediction, but evidence for a
greater \emph{framing-specific} effect in the author-relevant category is inconclusive at this
sample size.

On human--model alignment, the two components of H4 separate cleanly. Models track the
\emph{direction} of human shifts better than chance, but only weakly, which is what H4's word
``partially'' predicts. Humans show much larger shifts in \emph{magnitude}, by roughly 34
percentage points, and this survives standardising both sides to equal vote counts. A further
pattern, beyond magnitude, was not anticipated by the original hypotheses. It is reported
cautiously, since the two sides were not tested against each other directly: under framing, human
judgments become measurably less in agreement with one another, while the model ensemble shows no statistically
detectable dispersion effect in any category. Read as separate results rather than a demonstrated
contrast, they suggest that models redirect probability mass while human group judgments also
become more dispersed.

Two findings are reported as negative or near-negative results. Under multiple-comparison
correction, almost none of the individual per-emotion shifts remain significant: of 66 tests
spanning both sources, only one human effect (a gain in guilt on author-independent texts) and six
model effects survive, and the model effects form a coherent pattern of moving away from
\emph{joy} towards more specific positive emotions (\emph{pride}, \emph{relief}) rather than towards
threat-related ones. The exploratory comparison found no effect of prompt configuration on
framing-effect magnitude. One further result was unplanned: most of H1's support depends on the
choice of representation. The identical permutation test applied to the selected label rather than
the vote vector rejects the null in one of its six cases under a reverse-alphabetical tie-break
convention, and none under the alphabetical convention used throughout this thesis (only these two
otherwise-arbitrary conventions were tested). A majority-label version of this same study would
have detected a statistically significant human framing effect in one of those six tests at best,
or none under this thesis's own convention.

\section{Contribution}
\label{conc:contribution}

\noindent
Substantively, the thesis offers evidence of a real and measurable influence of perceived
authorship on the interpretation of emotion by both humans and language models. It adds a control
condition that isolates persona content from author presence, and it compares the two sides
directly under one design.

The methodological contribution may reach further still. Whether judgments are represented as
distributions or as labels is not a cosmetic difference; it determines whether the phenomenon can
be observed at all. A selected-label analysis of these same data would have recorded 38.7\% of
human and 41.4\% of model comparisons as showing no change when the distribution demonstrably
moved (under the alphabetical tie-break convention; crediting the label method for tracking the
complete set of tied top emotions instead lowers these to 25.8\% for humans and 38.7\% for models):
for models, more than twice the comparisons it would have detected, and for humans nearly
three-quarters as many as it would have detected. Even more directly, the same label-based
treatment finds \emph{no} significant human framing effect under the alphabetical tie-break
convention used as the primary analysis, and one of six under the reverse-alphabetical convention
also tested, where the distributional treatment finds three of six. Any study asking whether some
manipulation shifts subjective interpretation is exposed to this, and how large the blind spot is
only becomes clear once it is measured.

A second methodological result reaches just as far. Distributional distance between small vote
samples is not zero when there is no effect, and it scales with the dispersion of the underlying
distribution. Any ambiguity-defined grouping compares distances across groups that differ
systematically in dispersion, which confounds the effect of interest with the measurement's own
noise floor. An explicit per-item sampling null makes this visible, and in this study it changed
what could be concluded from H3.

\section{Closing Remarks}
\label{conc:closing}

\noindent
This thesis set out to challenge a widespread practice: collapsing multiple human judgments into
one gold label and treating the residue as noise \cite{Basile2020GoldStandard,
Aroyo2015CrowdTruth}. The practice is convenient and entrenched, but for a substantial fraction of
texts it is simply not well defined. Among the sampled texts (a sample stratified towards ambiguous
categories, not a corpus-wide rate), one text in five has no emotion commanding a majority of
readers, and 13.3\% have no unique most-popular answer at all. For the texts with tied top
choices, a pipeline reporting a single label depends on a tie-break; for texts with a unique
plurality but no majority, that label still conceals the disagreement.

The distributional view shows that the disagreement being thrown away is not simply noise. Its
overall magnitude and direction shift detectably when readers are told something about who wrote
the text, in humans and models alike, and the two share that shift weakly but measurably. The
claim concerns the aggregate shift, not every per-emotion pattern or category ordering examined
along the way; several of those remain inconclusive at this sample size. Emotion recognition
systems evaluated only against majority labels are scored on a target that discards variation this
thesis has shown to be both structured and responsive to context.

Should models reproduce human interpretive uncertainty, or is it enough that they give human-like
answers? That is a design question, not an empirical one. It cannot even be asked, though, while
the evaluation target is a single label, and the results reported here suggest that humans and
models differ in ways that matter.
